\paragraph{Immediate-finalizer control.}
A shrinking proposal-hit rate does not distinguish correct minorities that
were recovered from those that were simply lost. We therefore apply the same
\texttt{vote\_then\_judge\_tiebreak} finalizer independently to the proposals
saved at R0 through R5. This changes only the stage at which the finalizer is
applied; no debate generation is rerun. The analysis covers four societies and
five benchmarks, with equal weight assigned to each benchmark.

\begin{table}[H]
\centering
\caption{Immediate-finalizer and interaction control. Panel (a) compares the
R0 and R5 states by method. The final column is R5 Vote+Judge minus R0
Vote+Judge; confidence intervals resample questions within each benchmark
before recomputing the unweighted macro average. Panel (b) applies the same
finalizer at every saved stage and averages over all four societies and five
benchmarks. All accuracy, hit, gain, tie, and unanimity values are percentages.}
\label{tab:immediate_finalizer_control}
\vspace{0.3em}
\textbf{(a) Method-specific R0 and R5 comparison.}\\[0.2em]
\scriptsize
\resizebox{\textwidth}{!}{%
\begin{tabular}{lrrrrrrrr}
\toprule
Method & Vote R0 $\to$ R5 & Hit R0 $\to$ R5 & R0 V+J & R0 judge & R5 V+J & Additional R5 gain [95\% CI] & Tie R0 $\to$ R5 & Unanimous R0 $\to$ R5 \\
\midrule
Vanilla Debate & 75.65 $\to$ 77.99 & 83.03 $\to$ 78.98 & 76.03 & $+0.39$ & 77.89 & $+1.85\ [0.54,3.18]$ & 7.27 $\to$ 0.83 & 70.50 $\to$ 95.86 \\
Top-4 Debate & 76.25 $\to$ 78.16 & 84.11 $\to$ 79.27 & 76.84 & $+0.59$ & 78.42 & $+1.58\ [0.37,2.79]$ & 8.90 $\to$ 0.78 & 67.78 $\to$ 95.84 \\
AC-Greedy Debate & 76.84 $\to$ 78.41 & 84.93 $\to$ 79.60 & 77.59 & $+0.75$ & 78.51 & $+0.92\ [-0.39,2.27]$ & 9.58 $\to$ 0.91 & 64.61 $\to$ 95.45 \\
SC-Greedy Debate & 75.64 $\to$ 78.22 & 83.88 $\to$ 79.19 & 75.33 & $-0.31$ & 78.08 & $+2.74\ [1.40,4.11]$ & 9.56 $\to$ 0.63 & 63.07 $\to$ 96.11 \\
\bottomrule
\end{tabular}%
}

\vspace{0.7em}
\textbf{(b) Full stage trajectory.}\\[0.2em]
\small
\begin{tabular}{lrrrrrrr}
\toprule
Stage & Vote & Hit & Headroom & Vote+Judge & Gain over R0 vote & Tie & Unanimous \\
\midrule
R0 & 76.09 & 83.99 & 7.90 & 76.45 & $+0.36$ & 8.83 & 66.49 \\
R1 & 78.09 & 85.32 & 7.23 & 77.99 & $+1.89$ & 6.74 & 72.13 \\
R2 & 78.32 & 82.20 & 3.88 & 78.34 & $+2.25$ & 3.64 & 86.14 \\
R3 & 78.35 & 80.56 & 2.21 & 78.52 & $+2.42$ & 1.96 & 91.71 \\
R4 & 78.10 & 79.58 & 1.48 & 78.12 & $+2.02$ & 1.45 & 94.30 \\
R5 & 78.20 & 79.26 & 1.06 & 78.22 & $+2.13$ & 0.79 & 95.81 \\
\bottomrule
\end{tabular}
\end{table}

At R0, the finalizer contributes only 0.36 points on average. The R2 and R3
states reach 78.34\% and 78.52\% after the same finalizer, respectively 1.89
and 2.07 points above the R0 Vote+Judge result. This rise occurs while ties
fall from 8.83\% to 3.64\% at R2 and 1.96\% at R3, and unanimity rises from
66.49\% to 86.14\% and 91.71\%. The trajectory therefore shows rapid
interaction gains followed by a plateau, not a claim that the final round is
always optimal. At the method level, the additional R5 gain beyond the R0
readout ranges from 0.92 to 2.74 points; its interval excludes zero for
Vanilla, Top-4, and SC-Greedy.

\paragraph{Backbone-dependent conversion.}
We next ask whether initial headroom predicts final lift on OLMo as it does on
Qwen. For the pooled comparison, we use the four datasets and four society
constructions available for both backbones. The MMLU analysis pools the five
predefined Table~\ref{tab:main_results} subjects, and its Pearson interval is a
paired question-bootstrap interval.

\begin{table}[H]
\centering
\caption{Association between initial headroom $H=P_0-V_0$ and final lift
$G=F-V_0$ on the two backbones. The MMLU row reports a 95\% paired
question-bootstrap interval for Pearson $r$.}
\label{tab:backbone_headroom_conversion}
\small
\begin{tabular}{lcc}
\toprule
Scope & Qwen3-4B & OLMo3-7B \\
\midrule
Four shared datasets, pooled & $r=.770,\ \rho=.787$ & $r=.548,\ \rho=.648$ \\
Five MMLU subjects, pooled & $r=.568\ [.125,.787]$ & $r=.635\ [.219,.805]$ \\
\bottomrule
\end{tabular}
\end{table}

The positive OLMo associations show that proposal supply remains informative,
especially on MMLU. The failure mode is nevertheless dataset-dependent. On
OLMo MATH500, all four societies retain 3.33--6.33 points of initial headroom,
yet every final lift is negative, ranging from $-1.67$ to $-0.33$ points. This
localizes the failure to conversion rather than missing proposals: useful
answers are surfaced, but the backbone and readout do not reliably turn them
into a better final decision.
